% GENERATED FILE. Edit canonical lessons, then run npm run generate:books.
% canonical-source-hash: fnv1a64:46a36b037341ce81
% canonical-lessons: MR-C129-granthalay, MR-C129-upahargrh, MR-C129-tapal, MR-C129-imarat, MR-C129-nhanighar

\chapter{Library, Restaurant, Post, Building, Bathroom}
\label{ch:mr-library-restaurant-post-building-bathroom}
\hlchaptermodality{129}

\begin{chapteropening}
\textbf{By the end of this chapter:} \emph{I can say a library, a restaurant, the post, a building and a bathroom.}

Complete the last lesson of chapter 129: I can say a library, a restaurant, the post, a building and a bathroom.
\end{chapteropening}

\section[\texorpdfstring{\mr{ग्रंथालय} (granthālay)}{ग्रंथालय (granthālay)}]{\mr{ग्रंथालय} (granthālay) — a library}

\label{lesson:MR-C129-granthalay}



\begin{quote}
Before the new one: say the Marathi for an airport, then the Marathi for a hospital.

\emph{Read it:} the three signs aloud as one line — \textbf{\mr{उद्यान}, \mr{रुग्णालय}, \mr{विमानतळ}}
\end{quote}

\subsection*{You'll want to know: \mr{ग्रंथालय}}
\textbf{\mr{ग्रंथालय}} — \emph{granthālay} — \textquotedblleft{}a library\textquotedblright{}.

\begin{cousinweb}
\textbf{\mr{ग्रंथालय}} (\emph{granthālay}) is a compound of Sanskrit parts, a tatsama: \emph{grantha}, \textquotedblleft{}knot\textquotedblright{}, later \textquotedblleft{}book\textquotedblright{} (CDIAL 4350), and \emph{ālaya}, \textquotedblleft{}abode\textquotedblright{}.
\end{cousinweb}

\begin{grammarlens}[title={What you've built}]
One more word for this chapter.
\end{grammarlens}

\subsection*{Your turn}
\begin{itemize}
\raggedright
  \item \emph{Say it:} \emph{granthālay}
  \item \emph{Say it:} \emph{granthālay}, once more
  \item \emph{Say it:} say \emph{granthālay}
  \item \emph{Recall:} say the Marathi for an airport, then the Marathi for a hospital, then say \emph{granthālay} again
\end{itemize}

\begin{tcolorbox}[breakable,colback=teal!4,colframe=teal!35!black,title={Before you move on}]
What does \mr{ग्रंथालय} mean? (\textquotedblleft{}A library\textquotedblright{}.) Say it once more.
\end{tcolorbox}
\section[\texorpdfstring{\mr{उपहारगृह} (upahārgṛh)}{उपहारगृह (upahārgṛh)}]{\mr{उपहारगृह} (upahārgṛh) — a restaurant}

\label{lesson:MR-C129-upahargrh}



\begin{quote}
Before the new one: say the Marathi for a hospital, then the Marathi for a library.
\end{quote}

\subsection*{You'll want to know: \mr{उपहारगृह}}
\textbf{\mr{उपहारगृह}} — \emph{upahārgṛh} — \textquotedblleft{}a restaurant\textquotedblright{}.

\begin{cousinweb}
\textbf{\mr{उपहारगृह}} (\emph{upahārgṛh}) is a compound of Sanskrit parts, a tatsama: \emph{upahāra}, \textquotedblleft{}offering, gift of food\textquotedblright{}, and \emph{gṛha}, \textquotedblleft{}house\textquotedblright{}.
\end{cousinweb}

\begin{grammarlens}[title={What you've built}]
One more word for this chapter.
\end{grammarlens}

\subsection*{Your turn}
\begin{itemize}
\raggedright
  \item \emph{Say it:} \emph{upahārgṛh}
  \item \emph{Say it:} \emph{upahārgṛh}, once more
  \item \emph{Say it:} say \emph{upahārgṛh}
  \item \emph{Recall:} say the Marathi for a hospital, then the Marathi for a library, then say \emph{upahārgṛh} again
\end{itemize}

\begin{tcolorbox}[breakable,colback=teal!4,colframe=teal!35!black,title={Before you move on}]
What does \mr{उपहारगृह} mean? (\textquotedblleft{}A restaurant\textquotedblright{}.) Say it once more.
\end{tcolorbox}
\section[\texorpdfstring{\mr{टपाल} (ṭapāl)}{टपाल (ṭapāl)}]{\mr{टपाल} (ṭapāl) — the post, mail}

\label{lesson:MR-C129-tapal}



\begin{quote}
Before the new one: say the Marathi for a library, then the Marathi for a restaurant.
\end{quote}

\subsection*{You'll want to know: \mr{टपाल}}
\textbf{\mr{टपाल}} — \emph{ṭapāl} — \textquotedblleft{}the post, mail\textquotedblright{}.

\begin{cousinweb}
The origin of \textbf{\mr{टपाल}} (\emph{ṭapāl}) is not settled. Neither CDIAL nor DEDR, searched by form and by meaning, has an entry for it.
\end{cousinweb}

\begin{grammarlens}[title={What you've built}]
One more word for this chapter.
\end{grammarlens}

\subsection*{Your turn}
\begin{itemize}
\raggedright
  \item \emph{Say it:} \emph{ṭapāl}
  \item \emph{Say it:} \emph{ṭapāl}, once more
  \item \emph{Say it:} say \emph{ṭapāl}
  \item \emph{Recall:} say the Marathi for a library, then the Marathi for a restaurant, then say \emph{ṭapāl} again
\end{itemize}

\begin{tcolorbox}[breakable,colback=teal!4,colframe=teal!35!black,title={Before you move on}]
What does \mr{टपाल} mean? (\textquotedblleft{}The post, mail\textquotedblright{}.) Say it once more.
\end{tcolorbox}
\section[\texorpdfstring{\mr{इमारत} (imārat)}{इमारत (imārat)}]{\mr{इमारत} (imārat) — a building}

\label{lesson:MR-C129-imarat}



\begin{quote}
Before the new one: say the Marathi for a restaurant, then the Marathi for the post.
\end{quote}

\subsection*{You'll want to know: \mr{इमारत}}
\textbf{\mr{इमारत}} — \emph{imārat} — \textquotedblleft{}a building\textquotedblright{}.

\begin{cousinweb}
\textbf{\mr{इमारत}} (\emph{imārat}) is a loan from Arabic \emph{ʿimāra}, \textquotedblleft{}building, construction\textquotedblright{}, which reached Marathi through Persian. Its Arabic root is \emph{ʿ-m-r} (Lane). Persian and Hindustani say \emph{ʿimārat}.
\end{cousinweb}

\begin{grammarlens}[title={What you've built}]
One more word for this chapter.
\end{grammarlens}

\subsection*{Your turn}
\begin{itemize}
\raggedright
  \item \emph{Say it:} \emph{imārat}
  \item \emph{Say it:} \emph{imārat}, once more
  \item \emph{Say it:} say \emph{imārat}
  \item \emph{Recall:} say the Marathi for a restaurant, then the Marathi for the post, then say \emph{imārat} again
\end{itemize}

\begin{tcolorbox}[breakable,colback=teal!4,colframe=teal!35!black,title={Before you move on}]
What does \mr{इमारत} mean? (\textquotedblleft{}A building\textquotedblright{}.) Say it once more.
\end{tcolorbox}
\section[\texorpdfstring{\mr{न्हाणीघर} (nhāṇīghar)}{न्हाणीघर (nhāṇīghar)}]{\mr{न्हाणीघर} (nhāṇīghar) — a bathroom}

\label{lesson:MR-C129-nhanighar}



\begin{quote}
Before the new one: say the Marathi for the post, then the Marathi for a building.
\end{quote}

\subsection*{You'll want to know: \mr{न्हाणीघर}}
\textbf{\mr{न्हाणीघर}} — \emph{nhāṇīghar} — \textquotedblleft{}a bathroom\textquotedblright{}.

\begin{cousinweb}
\textbf{\mr{न्हाणीघर}} (\emph{nhāṇīghar}) is \textquotedblleft{}bathing house\textquotedblright{}: \emph{nhāṇī}, from Marathi \emph{nhāṇ}, \textquotedblleft{}bathing\textquotedblright{}, which CDIAL lists under Sanskrit \emph{snāna}, \textquotedblleft{}bathing\textquotedblright{} (CDIAL 13789), and \textbf{\mr{घर}}, \textquotedblleft{}house\textquotedblright{}.
\end{cousinweb}

\begin{grammarlens}[title={What you've built}]
One more word for this chapter.
\end{grammarlens}

\subsection*{Your turn}
\begin{itemize}
\raggedright
  \item \emph{Say it:} \emph{nhāṇīghar}
  \item \emph{Say it:} \emph{nhāṇīghar}, once more
  \item \emph{Say it:} say \emph{nhāṇīghar}
  \item \emph{Recall:} say the Marathi for the post, then the Marathi for a building, then say \emph{nhāṇīghar} again
\end{itemize}

\begin{tcolorbox}[breakable,colback=teal!4,colframe=teal!35!black,title={Before you move on}]
What does \mr{न्हाणीघर} mean? (\textquotedblleft{}A bathroom\textquotedblright{}.) Say it once more.
\end{tcolorbox}
